{\editcolor 
\vspace{-1mm}
\section{Limitations}
\label{sec:limitations}
\vspace{-1mm}

Our approach currently focuses on the connections between table segments and passages. 
{\feditcolor
In future work, we aim to extend our method to more general types of connections between nodes of arbitrary modalities, such as images and inter-table links. 
Concretely, we are actively exploring extensions to HELIOS’s graph construction and traversal mechanisms to incorporate image-based nodes alongside text and tables. 
At this stage, we convert image inputs into text representations using vision-language models (e.g., Qwen2.5-VL 7B), and treat the resulting summaries as nodes. 
Even with this simple transformation, preliminary experimental results indicate that HELIOS already outperforms existing multimodal SOTA retrievers. We plan to further develop this line of research and present it in a dedicated future publication.
Moreover, as shown in Section~\ref{sec:exp_ablation_study}, the effectiveness of our SLR module can be limited by LLM hallucinations. 
An interesting direction for future work is to explore self-evaluation techniques~\citep{wangself, zhang2024self}—for instance, computing confidence scores and re-checking low-confidence outputs—to further reduce any remaining hallucinations.
}
}